% =========================
% H) Reproducibility Statement (Exact vs Broad) + Artifact Plan
% =========================
\subsection{Reproducibility Statement (Exact vs Broad) + Artifact Plan}
\label{subsec:repro}

\paragraph{Exact reproducibility.}
We label experiments as \emph{exactly reproducible} when objective evaluations are deterministic:
either (i) the benchmark is tabular/deterministic, or (ii) the full training/evaluation pipeline can be made deterministic under our stack
by enforcing the determinism controls in \DeterminismConstraints{} (including fixed Python/NumPy/framework RNG seeds, fixed dataloader order,
and deterministic algorithm settings where supported).
In this regime, rerunning the same container with the same configurations and seed lists should reproduce results bitwise
(or as close to bitwise as the deterministic stack permits).
We therefore release the complete configurations, seed lists, and raw per-evaluation logs required to rerun these experiments exactly.

\paragraph{Broad reproducibility.}
For GPU-based training where nondeterminism may remain (e.g., non-deterministic kernels), experiments are \emph{broadly reproducible}:
reruns with the same container and seeds may yield small numeric differences, but results should match within statistical uncertainty
and preserve the reported conclusions (rankings, effect sizes, and confidence intervals).

\paragraph{Artifacts to release.}
We will release the following artifacts to support reproduction of all reported tables, figures, and statistical tests:
\begin{itemize}
    \item Source code for \ProposedMethodName{} and all runnable baselines (or pinned commits and wrappers for upstream implementations).
    \item A unified experiment harness and evaluator shared across methods (enforcing budgets, stopping, failure handling, and seeding).
    \item Machine-readable config files specifying benchmarks, methods, budgets, and seed lists.
    \item Environment specifications (Docker image / digest and/or Conda lockfile), plus \SoftwareVersions{}.
    \item Dataset version identifiers and preprocessing scripts, or immutable benchmark references where applicable.
    \item Raw per-evaluation logs (configurations, objective values, seeds, and failure codes) and aggregated results tables.
    \item Scripts to regenerate all figures and tables and to rerun the full statistical analysis pipeline.
    \item A run manifest recording benchmark id, method, budget, seed, worker count, git hash, and container hash.
    \item Artifact scope notes: \ArtifactReleaseScope{}.
\end{itemize}

\paragraph{Comparability and limitations.}
We explicitly delineate:
(i) which baselines are directly comparable under the same evaluation model,
(ii) which are not and why (see \BaselinesNotRunWhy{}),
and (iii) how this bounds claims (e.g., no claims about methods requiring different fidelities, priors, or search spaces).
